%% 中文摘要页
\begin{ChineseAbstract}
机器人在不确定环境中自主决策时，强化学习虽能提升策略性能，却难以保证控制过程不违反安全约束，且易受感知误差和训练—部署环境差异的影响；当任务扩展到多机器人编队，任务目标与安全约束、个体安全与群体结构之间的耦合使这一困难进一步加剧。针对这些问题，本文围绕单体安全保障与多机器人协同，提出约束流形策略学习方法、几何约束策略学习方法，并实现一套容器化仿真验证平台，分别从安全控制、编队策略与系统验证三个层面构成完整方案。

（1）针对单机器人自主决策中的安全保障问题，本文提出\textbf{约束流形策略学习方法}（CMPL, Constrained Manifold Policy Learning），由约束流形投影、安全可达性分析与数据驱动状态估计三部分构成统一框架。约束流形投影通过零空间映射将策略动作约束至安全流形的切空间，并结合奖励校准机制引导策略学习内生的安全行为；安全可达性分析基于历史轨迹数据预计算与环境相关的安全区域，在部署阶段无需在线重计算可达区域即可自适应调整安全裕度；数据驱动状态估计借助扩展卡尔曼滤波框架，由卷积神经网络学习观测噪声参数，使测量协方差能够随运动状态自适应调整，从而保证安全过滤器输入的可靠性。
实验结果表明，在目标到达任务消融实验中，CMPL 将平均安全代价由近端策略优化（PPO）基线方法的 50.58 降至 4.91，降幅达 90.3\%；在 Webots 仿真器中 E-puck 移动机器人零样本部署实验中，碰撞率由无安全过滤时的 40.0\% 降至在实验统计范围内未发生碰撞，位置误差由 0.051\,m 降至 0.037\,m，降幅 27.5\%。上述结果说明，CMPL 能够在显著降低约束违背风险的同时保持较好的迁移性与鲁棒性。

（2）针对多机器人编队导航中任务目标与安全约束相互耦合的情况，本文提出\textbf{几何约束策略学习方法}（GCPL, Geo-Constrained Policy Learning）。该方法以多智能体近端策略优化（MAPPO）为决策基础，在策略输出端引入几何约束过滤机制，对编队结构保持、机器人间安全距离以及避障要求进行统一处理，使安全性约束不再仅依赖奖励函数进行间接引导，而是在动作执行前完成显式修正。
实验结果显示，在线形编队任务中，GCPL 相较于 MAPPO 基线将安全代价由 105.78 降至 18.53，降幅 82.5\%；编队完成度由 10.00\% 提升至 71.00\%，成功率由 2.00\% 提升至 66.00\%。在圆形编队等复杂拓扑任务中，GCPL 将安全代价由 145.22 降至 34.40，降幅 76.3\%，成功率由 0.50\% 提升至 28.00\%。上述结果表明，GCPL 能够在保证安全性的同时显著提升任务完成效率与编队稳定性。

（3）本文设计并实现了\textbf{面向多机器人安全强化学习的容器化仿真验证平台}。该平台采用分层架构组织仿真环境、算法接口、评估模块与结果展示模块，通过容器化方式统一运行依赖，减少不同实验配置之间的环境差异；同时设置标准化评估流程，对不同算法在相同任务条件下进行公平比较。在此基础上，本文将 CMPL 与 GCPL 接入平台，并从安全约束满足与任务完成效果等维度进行综合实验分析，进一步验证了 CMPL 与 GCPL 在统一仿真环境下的端到端集成能力与综合性能。

\ChineseKeywords{约束流形；多智能体强化学习；多机器人编队；安全强化学习}
\end{ChineseAbstract}

%% 英文摘要页
\begin{EnglishAbstract}
    When a robot makes autonomous decisions in an uncertain environment, reinforcement learning improves task performance but cannot by itself keep the control process from violating safety constraints, and it is further degraded by perception errors and by the gap between training and deployment. Extending the task to multi-robot formations sharpens this difficulty, as task objectives now couple with safety constraints and individual safety couples with group structure. Addressing these issues, this thesis centers on single-robot safety assurance and multi-robot coordination: it proposes the Constrained Manifold Policy Learning method and the Geo-Constrained Policy Learning method, and implements a containerized simulation verification platform, together covering safety control, formation policy, and system-level verification.

    (1) For the safety assurance problem in single-robot autonomous decision-making, this thesis proposes the \textbf{Constrained Manifold Policy Learning} method (CMPL), a unified framework consisting of constrained-manifold projection, safe reachability analysis, and data-driven state estimation. The constrained-manifold projection maps policy actions onto the tangent space of the safety manifold via null-space projection and incorporates a reward calibration mechanism to drive the policy toward safe behavior. Safe reachability analysis pre-computes environment-dependent safety regions from historical trajectories so that the safety margin can be adapted at deployment without online recomputation of reachable regions. The data-driven state estimation employs an extended Kalman filter whose observation-noise parameters are learned by a convolutional neural network, allowing the measurement covariance to adapt with the motion state and ensuring reliable input to the safety filter. Experimental results show that on the goal-reaching ablation, CMPL reduces the average safety cost from 50.58 (PPO baseline) to 4.91, a reduction of 90.3\%; in zero-shot deployment of the E-puck robot in the Webots simulator, the collision rate decreases from 40.0\% (without safety filtering) to no observed collisions in the evaluation trials, and the position error from 0.051\,m to 0.037\,m, a 27.5\% reduction. These results demonstrate that the proposed method substantially reduces constraint-violation risk while preserving transferability and robustness.

    (2) For multi-robot formation navigation where task objectives and safety constraints are coupled, this thesis proposes the \textbf{Geo-Constrained Policy Learning} method (GCPL). Building on Multi-Agent Proximal Policy Optimization (MAPPO) as the decision backbone, GCPL introduces a geometric constraint filtering mechanism on the policy output that uniformly handles formation maintenance, inter-robot safe distance, and obstacle avoidance, so that safety constraints are no longer only indirectly steered by reward but are explicitly enforced before action execution.
    Experimental results show that on the line formation task, GCPL reduces the safety cost from 105.78 (MAPPO baseline) to 18.53 (an 82.5\% reduction), increases formation completion from 10.00\% to 71.00\%, and improves the success rate from 2.00\% to 66.00\%. On more complex topologies such as circle formation, GCPL reduces the safety cost from 145.22 to 34.40 (76.3\%) and increases the success rate from 0.50\% to 28.00\%. These results indicate that the proposed method significantly improves task efficiency and formation stability while preserving safety.
    
    (3) This thesis designs and implements a \textbf{containerized simulation verification platform for multi-robot safe reinforcement learning}. The platform adopts a layered architecture covering simulation environments, algorithm interfaces, evaluation modules, and result visualization, and uses containerization to unify runtime dependencies and minimize environment drift across experiment configurations. A standardized evaluation pipeline allows fair comparison of different algorithms under identical task conditions. On this platform, this thesis integrates the safety constraint mechanism with multi-robot cooperative policies and analyzes their behavior in terms of constraint satisfaction, task completion, and formation consistency, thereby verifying the overall effectiveness of the proposed methods at the system level.

    \EnglishKeywords{Constraint Manifold; Multi-agent Reinforcement Learning; Multi-robot Formation; Safe Reinforcement Learning}
\end{EnglishAbstract}
